\section{成化正德的朝廷、法律与文化连结}

本节以朝廷程序、法律语言、家庭与性别、城市文化和海外联系为问题线索。正史、实录、律例与奏议通常更容易保存官员、男性精英和国家分类；女性、家庭成员、工匠、城市居民与海外中介则常仅在墓志、契约、地方志、碑刻、航海记录或案件中留下片段。材料本身的偏向，应当成为叙事的一部分。

\subsection{内廷、司法与官僚程序}

\eventanchor{ML-1476-002}{明·1476（成化十二年）}
\paragraph{1476（成化十二年）：成化、弘治间中枢任免、厂卫与言路的本年材料标记皇权和官僚关系的调整。（材料核查重点：诏令或奏议的形成与发受链）} 此事把权力、法律或文化关系推进到可辨认的节点。账本的证据限度是来源导向的年度桥接条目：本年材料可定位相关政务线索；具体月日、发文机关、地域和执行结果仍须逐项回查，不能以制度文本或奏报替代实际效果。可资核查的材料为《宪宗实录》成化十二年条；《明史·本纪》及相关列传；诏令、奏疏或会典版本。它们能够说明制度如何表述自身、某种命令或交涉如何被记录，却未必能直接复原家庭内部关系、性别分工、城市日常或海外行动者的意图。事件发生的条件涉及继承、官僚协调或皇权运作的压力使问题进入朝廷程序；其后留下的制度与社会后果为它影响后续人事与政策议程；动机和派系标签不能由单一叙事裁定。桥接条目只标示可回查的年度问题与材料链；实录有中央选择性，崇祯期尤其需要奏疏、地方、朝鲜及满文材料交叉。；本条特别核对诏令或奏议的形成与发受链。

\eventanchor{ML-1476-003}{明·1476（成化十二年）}
\paragraph{1476（成化十二年）：成化、弘治间中枢任免、厂卫与言路的本年材料标记皇权和官僚关系的调整。（材料核查重点：报称信息的来源、抄传与行政分类）} 此事把权力、法律或文化关系推进到可辨认的节点。账本的证据限度是来源导向的年度桥接条目：本年材料可定位相关政务线索；具体月日、发文机关、地域和执行结果仍须逐项回查，不能以制度文本或奏报替代实际效果。可资核查的材料为《宪宗实录》成化十二年条；《明史·本纪》及相关列传；诏令、奏疏或会典版本。它们能够说明制度如何表述自身、某种命令或交涉如何被记录，却未必能直接复原家庭内部关系、性别分工、城市日常或海外行动者的意图。事件发生的条件涉及继承、官僚协调或皇权运作的压力使问题进入朝廷程序；其后留下的制度与社会后果为它影响后续人事与政策议程；动机和派系标签不能由单一叙事裁定。桥接条目只标示可回查的年度问题与材料链；实录有中央选择性，崇祯期尤其需要奏疏、地方、朝鲜及满文材料交叉。；本条特别核对报称信息的来源、抄传与行政分类。

\eventanchor{ML-1476-004}{明·1476（成化十二年）}
\paragraph{1476（成化十二年）：成化、弘治间中枢任免、厂卫与言路的本年材料标记皇权和官僚关系的调整。（材料核查重点：同年地方材料所见的执行范围）} 此事把权力、法律或文化关系推进到可辨认的节点。账本的证据限度是来源导向的年度桥接条目：本年材料可定位相关政务线索；具体月日、发文机关、地域和执行结果仍须逐项回查，不能以制度文本或奏报替代实际效果。可资核查的材料为《宪宗实录》成化十二年条；《明史·本纪》及相关列传；诏令、奏疏或会典版本。它们能够说明制度如何表述自身、某种命令或交涉如何被记录，却未必能直接复原家庭内部关系、性别分工、城市日常或海外行动者的意图。事件发生的条件涉及继承、官僚协调或皇权运作的压力使问题进入朝廷程序；其后留下的制度与社会后果为它影响后续人事与政策议程；动机和派系标签不能由单一叙事裁定。桥接条目只标示可回查的年度问题与材料链；实录有中央选择性，崇祯期尤其需要奏疏、地方、朝鲜及满文材料交叉。；本条特别核对同年地方材料所见的执行范围。

\eventanchor{MM-1476-YUNYANG-PREFECTURE}{明·1476（成化十二年）}
\paragraph{1476（成化十二年）：朝廷设郧阳府，以承接荆襄流民安置和山区行政} 此事把权力、法律或文化关系推进到可辨认的节点。账本的证据限度是设府时间与建置可靠；府治建立不能证明流民已固定定居，也不能代表整个山区的治理效果。可资核查的材料为《宪宗实录》成化十二年条；《明史·地理志》；《郧阳府志》建置沿革。它们能够说明制度如何表述自身、某种命令或交涉如何被记录，却未必能直接复原家庭内部关系、性别分工、城市日常或海外行动者的意图。事件发生的条件涉及荆襄招抚后，临时军政机构不足以处理户籍、赋役、司法与垦殖争端；其后留下的制度与社会后果为郧阳成为连接湖北、陕西、河南山地的行政节点；编户与土地控制仍在地方社会中持续协商。这也要求将正式名义、实际执行和当事人的处境分别讨论。

\eventanchor{ML-1477-001}{明·1477（成化十三年）}
\paragraph{1477（成化十三年）：本年灾荒、河工或赈恤的报奏材料可与地方志互校，报灾范围和实际损失应分别记录} 此事把权力、法律或文化关系推进到可辨认的节点。账本的证据限度是来源导向的年度桥接条目：本年材料可定位相关政务线索；具体月日、发文机关、地域和执行结果仍须逐项回查，不能以制度文本或奏报替代实际效果。可资核查的材料为《宪宗实录》成化十三年条；《明史·五行志》；受灾府县地方志与奏疏。它们能够说明制度如何表述自身、某种命令或交涉如何被记录，却未必能直接复原家庭内部关系、性别分工、城市日常或海外行动者的意图。事件发生的条件涉及自然风险与地方报告促成朝廷或地方的应对记录；其后留下的制度与社会后果为赈济、迁徙和损失的实际范围仍须以受灾地区材料分别核验。桥接条目只标示可回查的年度问题与材料链；实录有中央选择性，崇祯期尤其需要奏疏、地方、朝鲜及满文材料交叉。

\eventanchor{ML-1477-002}{明·1477（成化十三年）}
\paragraph{1477（成化十三年）：荆襄、广西、辽东或西北的兵事和安置文书构成本年地方秩序重组的线索。（材料核查重点：诏令或奏议的形成与发受链）} 此事把权力、法律或文化关系推进到可辨认的节点。账本的证据限度是来源导向的年度桥接条目：本年材料可定位相关政务线索；具体月日、发文机关、地域和执行结果仍须逐项回查，不能以制度文本或奏报替代实际效果。可资核查的材料为《宪宗实录》成化十三年条；《明史·兵志》及相关列传；边镇、地方志或对方材料。它们能够说明制度如何表述自身、某种命令或交涉如何被记录，却未必能直接复原家庭内部关系、性别分工、城市日常或海外行动者的意图。事件发生的条件涉及前期边防、地方武装或跨区域资源调动的压力推动了该行动；其后留下的制度与社会后果为它改变了后续调兵、补给或地方秩序的条件；战果和伤亡须分开核对。桥接条目只标示可回查的年度问题与材料链；实录有中央选择性，崇祯期尤其需要奏疏、地方、朝鲜及满文材料交叉。；本条特别核对诏令或奏议的形成与发受链。

\subsection{家庭、宗藩与性别规范}

\eventanchor{ML-1477-003}{明·1477（成化十三年）}
\paragraph{1477（成化十三年）：荆襄、广西、辽东或西北的兵事和安置文书构成本年地方秩序重组的线索。（材料核查重点：报称信息的来源、抄传与行政分类）} 此事把权力、法律或文化关系推进到可辨认的节点。账本的证据限度是来源导向的年度桥接条目：本年材料可定位相关政务线索；具体月日、发文机关、地域和执行结果仍须逐项回查，不能以制度文本或奏报替代实际效果。可资核查的材料为《宪宗实录》成化十三年条；《明史·兵志》及相关列传；边镇、地方志或对方材料。它们能够说明制度如何表述自身、某种命令或交涉如何被记录，却未必能直接复原家庭内部关系、性别分工、城市日常或海外行动者的意图。事件发生的条件涉及前期边防、地方武装或跨区域资源调动的压力推动了该行动；其后留下的制度与社会后果为它改变了后续调兵、补给或地方秩序的条件；战果和伤亡须分开核对。桥接条目只标示可回查的年度问题与材料链；实录有中央选择性，崇祯期尤其需要奏疏、地方、朝鲜及满文材料交叉。；本条特别核对报称信息的来源、抄传与行政分类。

\eventanchor{ML-1477-004}{明·1477（成化十三年）}
\paragraph{1477（成化十三年）：荆襄、广西、辽东或西北的兵事和安置文书构成本年地方秩序重组的线索。（材料核查重点：同年地方材料所见的执行范围）} 此事把权力、法律或文化关系推进到可辨认的节点。账本的证据限度是来源导向的年度桥接条目：本年材料可定位相关政务线索；具体月日、发文机关、地域和执行结果仍须逐项回查，不能以制度文本或奏报替代实际效果。可资核查的材料为《宪宗实录》成化十三年条；《明史·兵志》及相关列传；边镇、地方志或对方材料。它们能够说明制度如何表述自身、某种命令或交涉如何被记录，却未必能直接复原家庭内部关系、性别分工、城市日常或海外行动者的意图。事件发生的条件涉及前期边防、地方武装或跨区域资源调动的压力推动了该行动；其后留下的制度与社会后果为它改变了后续调兵、补给或地方秩序的条件；战果和伤亡须分开核对。桥接条目只标示可回查的年度问题与材料链；实录有中央选择性，崇祯期尤其需要奏疏、地方、朝鲜及满文材料交叉。；本条特别核对同年地方材料所见的执行范围。

\eventanchor{MM-1477-WESTERN-DEPOT}{明·1477（成化十三年）}
\paragraph{1477（成化十三年）：宪宗设置西厂，命宦官汪直领侦缉事务} 此事把权力、法律或文化关系推进到可辨认的节点。账本的证据限度是设厂和汪直职掌可靠；西厂案件范围、恐惧程度及其对全国的影响须区别实录、传记与后世笔记的修辞。可资核查的材料为《宪宗实录》成化十三年条；《明史·汪直传》；《明史·刑法志》。它们能够说明制度如何表述自身、某种命令或交涉如何被记录，却未必能直接复原家庭内部关系、性别分工、城市日常或海外行动者的意图。事件发生的条件涉及宪宗希望获得不经常规官僚渠道的案件与边情信息，汪直已在宫廷中获信任；其后留下的制度与社会后果为厂卫对官员和地方案件的介入扩大，引发文官反弹；西厂的存废成为内廷权力边界问题。这也要求将正式名义、实际执行和当事人的处境分别讨论。

\eventanchor{ML-1478-001}{明·1478（成化十四年）}
\paragraph{1478（成化十四年）：成化、弘治间中枢任免、厂卫与言路的本年材料标记皇权和官僚关系的调整。（材料核查重点：诏令或奏议的形成与发受链）} 此事把权力、法律或文化关系推进到可辨认的节点。账本的证据限度是来源导向的年度桥接条目：本年材料可定位相关政务线索；具体月日、发文机关、地域和执行结果仍须逐项回查，不能以制度文本或奏报替代实际效果。可资核查的材料为《宪宗实录》成化十四年条；《明史·本纪》及相关列传；诏令、奏疏或会典版本。它们能够说明制度如何表述自身、某种命令或交涉如何被记录，却未必能直接复原家庭内部关系、性别分工、城市日常或海外行动者的意图。事件发生的条件涉及继承、官僚协调或皇权运作的压力使问题进入朝廷程序；其后留下的制度与社会后果为它影响后续人事与政策议程；动机和派系标签不能由单一叙事裁定。桥接条目只标示可回查的年度问题与材料链；实录有中央选择性，崇祯期尤其需要奏疏、地方、朝鲜及满文材料交叉。；本条特别核对诏令或奏议的形成与发受链。

\eventanchor{ML-1478-002}{明·1478（成化十四年）}
\paragraph{1478（成化十四年）：成化、弘治间中枢任免、厂卫与言路的本年材料标记皇权和官僚关系的调整。（材料核查重点：报称信息的来源、抄传与行政分类）} 此事把权力、法律或文化关系推进到可辨认的节点。账本的证据限度是来源导向的年度桥接条目：本年材料可定位相关政务线索；具体月日、发文机关、地域和执行结果仍须逐项回查，不能以制度文本或奏报替代实际效果。可资核查的材料为《宪宗实录》成化十四年条；《明史·本纪》及相关列传；诏令、奏疏或会典版本。它们能够说明制度如何表述自身、某种命令或交涉如何被记录，却未必能直接复原家庭内部关系、性别分工、城市日常或海外行动者的意图。事件发生的条件涉及继承、官僚协调或皇权运作的压力使问题进入朝廷程序；其后留下的制度与社会后果为它影响后续人事与政策议程；动机和派系标签不能由单一叙事裁定。桥接条目只标示可回查的年度问题与材料链；实录有中央选择性，崇祯期尤其需要奏疏、地方、朝鲜及满文材料交叉。；本条特别核对报称信息的来源、抄传与行政分类。

\eventanchor{ML-1478-003}{明·1478（成化十四年）}
\paragraph{1478（成化十四年）：田土、赋役、盐课或漕运的本年政策记录提供财政执行的候选节点，需同地区册籍比照。（材料核查重点：诏令或奏议的形成与发受链）} 此事把权力、法律或文化关系推进到可辨认的节点。账本的证据限度是来源导向的年度桥接条目：本年材料可定位相关政务线索；具体月日、发文机关、地域和执行结果仍须逐项回查，不能以制度文本或奏报替代实际效果。可资核查的材料为《宪宗实录》成化十四年条；《明会典》相关条目；地方志、黄册/契约/河工材料。它们能够说明制度如何表述自身、某种命令或交涉如何被记录，却未必能直接复原家庭内部关系、性别分工、城市日常或海外行动者的意图。事件发生的条件涉及财政征解、交通工程或货币支付的压力使相关措施进入政务；其后留下的制度与社会后果为其法定安排不等于地方执行，后续须以地域材料追踪负担与成效。桥接条目只标示可回查的年度问题与材料链；实录有中央选择性，崇祯期尤其需要奏疏、地方、朝鲜及满文材料交叉。；本条特别核对诏令或奏议的形成与发受链。

\subsection{城市、消费与文化空间}

\eventanchor{ML-1478-004}{明·1478（成化十四年）}
\paragraph{1478（成化十四年）：田土、赋役、盐课或漕运的本年政策记录提供财政执行的候选节点，需同地区册籍比照。（材料核查重点：报称信息的来源、抄传与行政分类）} 此事把权力、法律或文化关系推进到可辨认的节点。账本的证据限度是来源导向的年度桥接条目：本年材料可定位相关政务线索；具体月日、发文机关、地域和执行结果仍须逐项回查，不能以制度文本或奏报替代实际效果。可资核查的材料为《宪宗实录》成化十四年条；《明会典》相关条目；地方志、黄册/契约/河工材料。它们能够说明制度如何表述自身、某种命令或交涉如何被记录，却未必能直接复原家庭内部关系、性别分工、城市日常或海外行动者的意图。事件发生的条件涉及财政征解、交通工程或货币支付的压力使相关措施进入政务；其后留下的制度与社会后果为其法定安排不等于地方执行，后续须以地域材料追踪负担与成效。桥接条目只标示可回查的年度问题与材料链；实录有中央选择性，崇祯期尤其需要奏疏、地方、朝鲜及满文材料交叉。；本条特别核对报称信息的来源、抄传与行政分类。

\eventanchor{ML-1478-005}{明·1478（成化十四年）}
\paragraph{1478（成化十四年）：田土、赋役、盐课或漕运的本年政策记录提供财政执行的候选节点，需同地区册籍比照。（材料核查重点：同年地方材料所见的执行范围）} 此事把权力、法律或文化关系推进到可辨认的节点。账本的证据限度是来源导向的年度桥接条目：本年材料可定位相关政务线索；具体月日、发文机关、地域和执行结果仍须逐项回查，不能以制度文本或奏报替代实际效果。可资核查的材料为《宪宗实录》成化十四年条；《明会典》相关条目；地方志、黄册/契约/河工材料。它们能够说明制度如何表述自身、某种命令或交涉如何被记录，却未必能直接复原家庭内部关系、性别分工、城市日常或海外行动者的意图。事件发生的条件涉及财政征解、交通工程或货币支付的压力使相关措施进入政务；其后留下的制度与社会后果为其法定安排不等于地方执行，后续须以地域材料追踪负担与成效。桥接条目只标示可回查的年度问题与材料链；实录有中央选择性，崇祯期尤其需要奏疏、地方、朝鲜及满文材料交叉。；本条特别核对同年地方材料所见的执行范围。

\eventanchor{ML-1479-001}{明·1479（成化十五年）}
\paragraph{1479（成化十五年）：四川苗族起义} 此事把权力、法律或文化关系推进到可辨认的节点。账本的证据限度是编年候选的年序可由官修与后出材料交叉；具体月日、参与者、战果或数字必须回查所列材料，不能将单一叙事当作定论。可资核查的材料为《宪宗实录》成化十五年条；《明史·本纪》及相关列传；诏令、奏疏或会典版本。它们能够说明制度如何表述自身、某种命令或交涉如何被记录，却未必能直接复原家庭内部关系、性别分工、城市日常或海外行动者的意图。事件发生的条件涉及继承、官僚协调或皇权运作的压力使问题进入朝廷程序；其后留下的制度与社会后果为它影响后续人事与政策议程；动机和派系标签不能由单一叙事裁定。译写候选以编年材料定位；实录记载反映朝廷选择，崇祯期须另以奏疏、地方及敌对政权材料交叉。

\eventanchor{ML-1479-002}{明·1479（成化十五年）四月28日}
\paragraph{1479（成化十五年）四月28日：越南黎朝皇帝黎浩向中国宫廷赠送金银器皿、当地丝绸制品} 此事把权力、法律或文化关系推进到可辨认的节点。账本的证据限度是编年候选的年序可由官修与后出材料交叉；具体月日、参与者、战果或数字必须回查所列材料，不能将单一叙事当作定论。可资核查的材料为《宪宗实录》成化十五年条；《明史·外国传》；《朝鲜王朝实录》、琉球《历代宝案》或海外档案。它们能够说明制度如何表述自身、某种命令或交涉如何被记录，却未必能直接复原家庭内部关系、性别分工、城市日常或海外行动者的意图。事件发生的条件涉及朝贡名义、边境安全与港口贸易的利益使交涉持续发生；其后留下的制度与社会后果为它为后续册封、互市或海防安排留下文书与跨政权比较线索。译写候选以编年材料定位；实录记载反映朝廷选择，崇祯期须另以奏疏、地方及敌对政权材料交叉。

\eventanchor{ML-1479-003}{明·1479（成化十五年）}
\paragraph{1479（成化十五年）：哈密、吐鲁番及周边朝贡关系的本年交涉显示东天山政治流动，册封不等于稳定控制。（材料核查重点：诏令或奏议的形成与发受链）} 此事把权力、法律或文化关系推进到可辨认的节点。账本的证据限度是来源导向的年度桥接条目：本年材料可定位相关政务线索；具体月日、发文机关、地域和执行结果仍须逐项回查，不能以制度文本或奏报替代实际效果。可资核查的材料为《宪宗实录》成化十五年条；《明史·外国传》；《朝鲜王朝实录》、琉球《历代宝案》或海外档案。它们能够说明制度如何表述自身、某种命令或交涉如何被记录，却未必能直接复原家庭内部关系、性别分工、城市日常或海外行动者的意图。事件发生的条件涉及朝贡名义、边境安全与港口贸易的利益使交涉持续发生；其后留下的制度与社会后果为它为后续册封、互市或海防安排留下文书与跨政权比较线索。桥接条目只标示可回查的年度问题与材料链；实录有中央选择性，崇祯期尤其需要奏疏、地方、朝鲜及满文材料交叉。；本条特别核对诏令或奏议的形成与发受链。

\eventanchor{ML-1479-004}{明·1479（成化十五年）}
\paragraph{1479（成化十五年）：哈密、吐鲁番及周边朝贡关系的本年交涉显示东天山政治流动，册封不等于稳定控制。（材料核查重点：报称信息的来源、抄传与行政分类）} 此事把权力、法律或文化关系推进到可辨认的节点。账本的证据限度是来源导向的年度桥接条目：本年材料可定位相关政务线索；具体月日、发文机关、地域和执行结果仍须逐项回查，不能以制度文本或奏报替代实际效果。可资核查的材料为《宪宗实录》成化十五年条；《明史·外国传》；《朝鲜王朝实录》、琉球《历代宝案》或海外档案。它们能够说明制度如何表述自身、某种命令或交涉如何被记录，却未必能直接复原家庭内部关系、性别分工、城市日常或海外行动者的意图。事件发生的条件涉及朝贡名义、边境安全与港口贸易的利益使交涉持续发生；其后留下的制度与社会后果为它为后续册封、互市或海防安排留下文书与跨政权比较线索。桥接条目只标示可回查的年度问题与材料链；实录有中央选择性，崇祯期尤其需要奏疏、地方、朝鲜及满文材料交叉。；本条特别核对报称信息的来源、抄传与行政分类。

\subsection{海外联系与外交分类}

\eventanchor{ML-1479-005}{明·1479（成化十五年）}
\paragraph{1479（成化十五年）：哈密、吐鲁番及周边朝贡关系的本年交涉显示东天山政治流动，册封不等于稳定控制。（材料核查重点：同年地方材料所见的执行范围）} 此事把权力、法律或文化关系推进到可辨认的节点。账本的证据限度是来源导向的年度桥接条目：本年材料可定位相关政务线索；具体月日、发文机关、地域和执行结果仍须逐项回查，不能以制度文本或奏报替代实际效果。可资核查的材料为《宪宗实录》成化十五年条；《明史·外国传》；《朝鲜王朝实录》、琉球《历代宝案》或海外档案。它们能够说明制度如何表述自身、某种命令或交涉如何被记录，却未必能直接复原家庭内部关系、性别分工、城市日常或海外行动者的意图。事件发生的条件涉及朝贡名义、边境安全与港口贸易的利益使交涉持续发生；其后留下的制度与社会后果为它为后续册封、互市或海防安排留下文书与跨政权比较线索。桥接条目只标示可回查的年度问题与材料链；实录有中央选择性，崇祯期尤其需要奏疏、地方、朝鲜及满文材料交叉。；本条特别核对同年地方材料所见的执行范围。

\eventanchor{ML-1480-001}{明·1480（成化十六年）}
\paragraph{1480（成化十六年）：荆襄、广西、辽东或西北的兵事和安置文书构成本年地方秩序重组的线索。（材料核查重点：诏令或奏议的形成与发受链）} 此事把权力、法律或文化关系推进到可辨认的节点。账本的证据限度是来源导向的年度桥接条目：本年材料可定位相关政务线索；具体月日、发文机关、地域和执行结果仍须逐项回查，不能以制度文本或奏报替代实际效果。可资核查的材料为《宪宗实录》成化十六年条；《明史·兵志》及相关列传；边镇、地方志或对方材料。它们能够说明制度如何表述自身、某种命令或交涉如何被记录，却未必能直接复原家庭内部关系、性别分工、城市日常或海外行动者的意图。事件发生的条件涉及前期边防、地方武装或跨区域资源调动的压力推动了该行动；其后留下的制度与社会后果为它改变了后续调兵、补给或地方秩序的条件；战果和伤亡须分开核对。桥接条目只标示可回查的年度问题与材料链；实录有中央选择性，崇祯期尤其需要奏疏、地方、朝鲜及满文材料交叉。；本条特别核对诏令或奏议的形成与发受链。

\eventanchor{ML-1480-002}{明·1480（成化十六年）}
\paragraph{1480（成化十六年）：荆襄、广西、辽东或西北的兵事和安置文书构成本年地方秩序重组的线索。（材料核查重点：报称信息的来源、抄传与行政分类）} 此事把权力、法律或文化关系推进到可辨认的节点。账本的证据限度是来源导向的年度桥接条目：本年材料可定位相关政务线索；具体月日、发文机关、地域和执行结果仍须逐项回查，不能以制度文本或奏报替代实际效果。可资核查的材料为《宪宗实录》成化十六年条；《明史·兵志》及相关列传；边镇、地方志或对方材料。它们能够说明制度如何表述自身、某种命令或交涉如何被记录，却未必能直接复原家庭内部关系、性别分工、城市日常或海外行动者的意图。事件发生的条件涉及前期边防、地方武装或跨区域资源调动的压力推动了该行动；其后留下的制度与社会后果为它改变了后续调兵、补给或地方秩序的条件；战果和伤亡须分开核对。桥接条目只标示可回查的年度问题与材料链；实录有中央选择性，崇祯期尤其需要奏疏、地方、朝鲜及满文材料交叉。；本条特别核对报称信息的来源、抄传与行政分类。

\eventanchor{ML-1480-003}{明·1480（成化十六年）}
\paragraph{1480（成化十六年）：学校、考试、刻书或礼制材料在本年留下文化制度线索，精英文本不能代表全体社会。（材料核查重点：诏令或奏议的形成与发受链）} 此事把权力、法律或文化关系推进到可辨认的节点。账本的证据限度是来源导向的年度桥接条目：本年材料可定位相关政务线索；具体月日、发文机关、地域和执行结果仍须逐项回查，不能以制度文本或奏报替代实际效果。可资核查的材料为《宪宗实录》成化十六年条；《明史·选举志》或《艺文志》；文集、碑刻或书目材料。它们能够说明制度如何表述自身、某种命令或交涉如何被记录，却未必能直接复原家庭内部关系、性别分工、城市日常或海外行动者的意图。事件发生的条件涉及制度、教育或知识传播的需要推动了相关行动或文本形成；其后留下的制度与社会后果为它提供制度和传播史的锚点，不能据此推断社会整体接受。桥接条目只标示可回查的年度问题与材料链；实录有中央选择性，崇祯期尤其需要奏疏、地方、朝鲜及满文材料交叉。；本条特别核对诏令或奏议的形成与发受链。

\eventanchor{ML-1480-004}{明·1480（成化十六年）}
\paragraph{1480（成化十六年）：学校、考试、刻书或礼制材料在本年留下文化制度线索，精英文本不能代表全体社会。（材料核查重点：报称信息的来源、抄传与行政分类）} 此事把权力、法律或文化关系推进到可辨认的节点。账本的证据限度是来源导向的年度桥接条目：本年材料可定位相关政务线索；具体月日、发文机关、地域和执行结果仍须逐项回查，不能以制度文本或奏报替代实际效果。可资核查的材料为《宪宗实录》成化十六年条；《明史·选举志》或《艺文志》；文集、碑刻或书目材料。它们能够说明制度如何表述自身、某种命令或交涉如何被记录，却未必能直接复原家庭内部关系、性别分工、城市日常或海外行动者的意图。事件发生的条件涉及制度、教育或知识传播的需要推动了相关行动或文本形成；其后留下的制度与社会后果为它提供制度和传播史的锚点，不能据此推断社会整体接受。桥接条目只标示可回查的年度问题与材料链；实录有中央选择性，崇祯期尤其需要奏疏、地方、朝鲜及满文材料交叉。；本条特别核对报称信息的来源、抄传与行政分类。

\eventanchor{ML-1481-001}{明·1481（成化十七年）}
\paragraph{1481（成化十七年）：田土、赋役、盐课或漕运的本年政策记录提供财政执行的候选节点，需同地区册籍比照。（材料核查重点：诏令或奏议的形成与发受链）} 此事把权力、法律或文化关系推进到可辨认的节点。账本的证据限度是来源导向的年度桥接条目：本年材料可定位相关政务线索；具体月日、发文机关、地域和执行结果仍须逐项回查，不能以制度文本或奏报替代实际效果。可资核查的材料为《宪宗实录》成化十七年条；《明会典》相关条目；地方志、黄册/契约/河工材料。它们能够说明制度如何表述自身、某种命令或交涉如何被记录，却未必能直接复原家庭内部关系、性别分工、城市日常或海外行动者的意图。事件发生的条件涉及财政征解、交通工程或货币支付的压力使相关措施进入政务；其后留下的制度与社会后果为其法定安排不等于地方执行，后续须以地域材料追踪负担与成效。桥接条目只标示可回查的年度问题与材料链；实录有中央选择性，崇祯期尤其需要奏疏、地方、朝鲜及满文材料交叉。；本条特别核对诏令或奏议的形成与发受链。
